\section{The Tutte polynomial and planar duality}\label{fl:tutte}

Let $G=(V,E)$ now be an undirected multigraph with indexed edges, including
loops. Put $n=|V|$, $m=|E|$, let $c(A)$ count components of $(V,A)$
including isolated vertices, and write $r(A)=n-c(A)$, $r=r(E)$. A spanning
forest of $(V,A)$ that connects each of its components uses $r(A)$ edges.
This follows by growing a tree one vertex at a time in each component.
In particular $r(A)\le |A|$. Define
\[
 T_G(x,y)=\sum_{A\subseteq E}(x-1)^{r-r(A)}(y-1)^{|A|-r(A)}.
\]
All exponents are nonnegative. Coefficients form a signed rectangular
array with coordinate bounds $n,m$. Expanding each summand by the
binomial rule gives an absolute bound $2^{n+2m}$ on each coefficient,
since there are $2^m$ summands, each having sum of absolute coefficients
at most $2^{n+m}$. Thus all polynomial identities below are arithmetic
statements about finite arrays.

\begin{theorem}[Deletion, contraction and special values]\leavevmode\label{fl:tutterec}
The edgeless graph has $T_G=1$. For a loop $e$, $T_G=yT_{G-e}$.
For a bridge $e$, $T_G=xT_{G/e}$. For every other edge,
$T_G=T_{G-e}+T_{G/e}$. In particular all the coefficients of $T_G$ in
$x,y$ are nonnegative, it is multiplicative under disjoint union, and
\[
\begin{array}{c|l}
 (x,y)&T_G(x,y)\ \text{counts}\\ \hline
 (1,1)&\text{spanning forests connecting each component of }G\\
 (2,1)&\text{all acyclic spanning edge subsets}\\
 (1,2)&\text{spanning edge subsets with the same components as }G\\
 (2,2)&\text{all }2^m\text{ edge subsets}.
\end{array}
\]
For a connected nonempty graph the first entry is its spanning-tree count.
\end{theorem}
\begin{proof}
For $B\subseteq E\setminus\{e\}$ pair the terms $B,B\cup\{e\}$.
A loop changes only the edge count, so the second term is $y-1$
times the first. For a nonloop, contracting $e$ identifies exactly
its two ends and retains the indices of all other edges, keeping any
new loops. Hence $r_{G/e}(B)=r_G(B\cup\{e\})-1$ and
$r_{G/e}=r-1$. The terms containing $e$ give $T_{G/e}$. If $e$ is
not a bridge, deletion leaves $r$ fixed and terms omitting $e$ give
$T_{G-e}$. If it is a bridge, its ends are in distinct components of
every $(V,B)$. Then $r(B\cup\{e\})=r(B)+1$ and the term for $B$ is
$x-1$ times the term for $B\cup\{e\}$. This gives the bridge rule.

The recurrence starts at one and uses only addition and multiplication
by a variable, proving coefficient nonnegativity. An edge-index order
bounds every recursion branch by $m$ deletions or contractions. A binary
tree of depth $m$ can be evaluated bottom-up in a bounded table.
For disjoint unions all ranks, sizes and exponent differences add
over the two factors. Finally, $|A|-r(A)=0$ exactly when $A$ is a
forest: a spanning forest inside $A$ has $r(A)$ edges and every extra
edge closes a cycle, possibly a loop or a parallel pair. The other
exponent vanishes exactly when $A$ connects every original component.
Substituting the four pairs, with $0^0=1$, proves the table.
\end{proof}

\begin{theorem}[Chromatic and flow evaluations]\leavevmode\label{fl:tutteflow}
For a natural $q$, let $P_G(q)$ count proper $q$-colourings, with a loop
forbidding every colouring. Give every edge an arbitrary direction.
For a finite abelian group of order $q\ge1$, a circulation assigns a
group element to each edge so that the sum entering each vertex equals
the sum leaving it. Let $F_G(q)$ count circulations nonzero on every
edge. Then
\[
 P_G(q)=(-1)^r q^{c(E)}T_G(1-q,0),\qquad
 F_G(q)=(-1)^{m-r}T_G(0,1-q).
\]
In particular the second count depends only on the group order and
is independent of the orientation.
\end{theorem}
\begin{proof}
The alternating sum $\sum_{A\subseteq E}(-1)^{|A|}q^{c(A)}$
counts proper colourings. A word constant on the ends of every edge
in $A$ has one free colour per component. Interchanging the finite
sums, an improper word cancels by toggling its first monochromatic
edge, and a proper one contributes one. This includes loops.
Substitution in $T$ gives this same sum, since $c(E)+r-r(A)=c(A)$
and the sign exponent is congruent to $|A|$ modulo two.

The number of circulations supported in $A$ (zero outside $A$) is
$q^{|A|-r(A)}$. Indeed choose a spanning forest inside $A$ and set
the values of the remaining $|A|-r(A)$ edges freely. In each
nontrivial tree take the first leaf. Its unique remaining tree edge
has a unique value balancing the leaf, obtained by moving all known
terms to the opposite side and using group inversion when needed.
Remove the leaf, retaining the contributions of the now fixed edge at
its other end. After at most $n$ steps only the roots remain. In each
component the sum of all vertex imbalances was zero for every full
assignment, by cancellation of each oriented edge. Since the nonroot
imbalances vanish, the root has balance zero too. This construction
works with any finite addition table satisfying the abelian group laws.

Now sum these counts with weight $(-1)^{m-|A|}$. A fixed circulation
is counted once for every superset of its support. If some edge
carries zero, toggling the first such edge pairs opposite signs.
A nowhere-zero circulation contributes one. The resulting sum equals
the second evaluation, because its sign on $A$ is
$(-1)^{m-r+r-r(A)+|A|-r(A)}=(-1)^{m-|A|}$.
\end{proof}

\begin{theorem}[Spherical duality]\leavevmode\label{fl:tuttedual}
Take a connected graph with a supplied orientable rotation system of
Euler characteristic two. Its face dual $G^*$ has a vertex per face
and a corresponding edge crossing each primal edge. With these indexed
edges identified,
\[
 T_{G^*}(x,y)=T_G(y,x),\qquad
 P_{G^*}(q)=q F_G(q)\quad(q\ge1).
\]
\end{theorem}
\begin{proof}
Write $f$ for the number of faces. In the cyclic order of half-edges at
each vertex, let $\sigma$ send each half-edge to the next one and
$\iota$ exchange the halves of each edge. Face walks are the cycles of
$\sigma\iota$ (reversing the rotation convention gives an equivalent
proof). The action generated by $\sigma,\iota$ is transitive on the
half-edges of a connected graph with an edge: follow a vertex path,
rotating at each intermediate vertex to its desired exiting half-edge.
The same action is generated by $\sigma\iota,\iota$, so the face dual
is connected as well. When $m=0$ there is one vertex and one face,
and both assertions follow directly.

In binary edge masks, the cuts of the dual are even subgraphs of the
primal graph. Indeed the cut of one dual vertex is the parity mask of
edges in that face's closed boundary walk. Every vertex there has even
degree modulo two. A general cut is the symmetric difference of these
single-face cuts and stays even. The connected dual has $2^{f-1}$ cuts:
a vertex mask and its complement give the same cut, and equality of
cuts makes membership in their symmetric difference constant along
every edge, hence constant throughout. The primal has $2^{m-r}$ even
subgraphs by the circulation count for the group of order two with
all edges allowed. Since $n-m+f=2$ and $r=n-1$, the counts are equal.
Membership in the two families is decidable on the same $m$-bit
universe, so inclusion and equality of finite counts give equality
of the families.

Fix $A\subseteq E$, and write $B=E\setminus A$. Counting only masks
supported in $A$ gives $2^{|A|-r(A)}$ on the primal side. On the dual
side a defining vertex mask must be constant on each component of
$(V^*,B)$. There are $2^{c^*(B)}$ choices and each cut is counted twice.
Consequently
\[
 |A|-r(A)=c^*(B)-1=r^*(E)-r^*(B).
\]
Since $r^*(E)=m-r$, rearrangement also gives
$r-r(A)=|B|-r^*(B)$. Complementing edge masks therefore exchanges
the two exponents term by term in the Tutte sum. Finally apply the two
preceding evaluations, using $c^*(E)=1$ and $r^*(E)=m-r$. Every face
scan uses at most $2m$ half-edges and all comparisons of cut and cycle
families are scans of coded finite masks.
\end{proof}

\begin{corollary}[Tait's formulation for a cubic spherical map]\leavevmode\label{thm:tait}
Suppose $G$ is a nonempty connected loopless cubic graph with a spherical
rotation. Its faces admit four colours, with different colours across every
edge, exactly when its edges admit a proper three-colouring. Both conversions
are primitive recursive.
\end{corollary}
\begin{proof}
Take the additive group $K=(\mathbb Z/2)^2$, with three nonzero elements.
A nowhere-zero $K$-flow on a cubic graph assigns three nonzero labels at each
vertex summing to zero. Orientation is immaterial in this group. Three such
labels sum to zero exactly when they are the three different nonzero elements.
Indeed two equal labels cancel, leaving a nonzero label, while the sum of all
three nonzero elements is zero. Thus these flows are exactly proper
edge three-colourings, using those labels as colours. A face colouring is a
vertex colouring of the dual, with a dual loop making proper colouring
impossible. Theorem~\ref{fl:tuttedual} gives
$P_{G^*}(4)=4F_G(4)$. This is an identity between counts of two explicit
finite tables, so given a witness on either side one may scan the other
table (at most $4^f$ or $3^m$ words) for a witness there. More directly,
from a proper face colouring, label each edge by the sum of its two face
colours. It is nonzero and each incident face sector contributes twice
around a vertex, so the three labels sum to zero. Both arguments cover the
case of a bridge, which borders the same face on both sides and is forced
to carry zero in a circulation.
\end{proof}
